O controle de ventilação em ambientes internos é essencial para garantir conforto térmico, qualidade do ar e eficiência energética. Métodos tradicionais podem ser limitados diante das incertezas e variabilidades presentes em ambientes reais. Nesse contexto, sistemas fuzzy, especialmente o modelo Mamdani~\cite{mamdani1975experiment}, destacam-se por lidar com informações imprecisas e incorporar conhecimento linguístico.

Este trabalho apresenta o desenvolvimento de um sistema de controle fuzzy Mamdani para ventilação, utilizando como variáveis de entrada a temperatura, a umidade e o número de pessoas no ambiente. O objetivo é demonstrar como a lógica fuzzy pode ser aplicada para tomada de decisão automática, promovendo conforto térmico e eficiência energética.

As variáveis de entrada e saída foram particionadas em três conjuntos fuzzy (baixa, média e alta para as entradas; fraca, moderada e forte para a saída), e a decisão é tomada com base em uma base de regras fuzzy fundamentada em conhecimentos de conforto térmico.

A metodologia adotada compreende a definição dos universos de discurso, a construção das funções de pertinência triangulares, a elaboração da base de regras fuzzy, a aplicação dos operadores de inferência Mamdani, a agregação das saídas das regras e a defuzzificação utilizando diferentes técnicas. Os resultados são apresentados em formato gráfico e quantitativo, permitindo análise crítica e discussão sobre o desempenho do sistema fuzzy para o problema proposto.